%%%%%%%%%%%%%%%%%%%Macros

%%%%%%%%%%%%%%%%%%% Layout related %%%%%%%%%%%%%%%%

% Clear double page without any header or footer at end of chapters
\newcommand{\emptydoublepage}{\clearpage\thispagestyle{empty}\cleardoublepage}

% Global scaling parameters
\newlength{\rulethickness}
\setlength{\rulethickness}{0.4pt}

%%%%%%%%%%%%%%%% Colors %%%%%%%%%%%%%%%%%%%%%%%%%%%
\colorlet{goodArea}{green!50}
\colorlet{badArea}{red!50}


%%%%%%%%%%%%%%%%%% TikZ %%%%%%%%%%%%%%%%%%%%%%%%%%%%%%%
\tikzset{
state/.style={draw,circle, inner sep=0pt, minimum size=20pt},
action/.style={draw, circle,  inner sep=0.75pt, fill},
probabilistic/.style={->},
markovian/.style={dashed, ->},
}

%%%%%%%%%%%%%%%% algorithmicx %%%%%%%%%%%%%%%%%%%%%%%%
\algnewcommand\algorithmicinput{\textbf{Input:}}
\algnewcommand\INPUT{\item[\algorithmicinput]}
\algnewcommand\algorithmicoutput{\textbf{Output:}}
\algnewcommand\OUTPUT{\item[\algorithmicoutput]}
\newcommand{\setLineNumber}[1]{\makeatletter\setcounter{ALG@line}{#1}\addtocounter{ALG@line}{-1}\makeatother}
\newcommand{\StateNoLineBreak}{\!;\ \ }
\newcommand{\StateSmallEmptyLine}{\Statex \vspace{-7pt}}
\renewcommand{\algorithmiccomment}[1]{/\!/ \textit{#1}}


%%%%%%%%%%%%%%%%%% Misc Math stuff $$$$$$$$$$$$$$$$$$$$


\newcounter{claimcounter}
\newcommand{\tagclaim}{\stepcounter{claimcounter}\tag{Claim~\arabic{claimcounter}}}

% xdashrightarrow
% http://tex.stackexchange.com/questions/136772/dotted-or-dashed-xrightarrow
\makeatletter
\newcommand*{\da@rightarrow}{\mathchar"0\hexnumber@\symAMSa 4B }
\newcommand*{\da@leftarrow}{\mathchar"0\hexnumber@\symAMSa 4C }
\newcommand*{\xdashrightarrow}[2][]{%
	\mathrel{%
		\mathpalette{\da@xarrow{#1}{#2}{}\da@rightarrow{\,}{}}{}%
	}%
}
\newcommand{\xdashleftarrow}[2][]{%
	\mathrel{%
		\mathpalette{\da@xarrow{#1}{#2}\da@leftarrow{}{}{\,}}{}%
	}%
}
\newcommand*{\da@xarrow}[7]{%
	% #1: below
	% #2: above
	% #3: arrow left
	% #4: arrow right
	% #5: space left 
	% #6: space right
	% #7: math style 
	\sbox0{$\ifx#7\scriptstyle\scriptscriptstyle\else\scriptstyle\fi#5#1#6\m@th$}%
	\sbox2{$\ifx#7\scriptstyle\scriptscriptstyle\else\scriptstyle\fi#5#2#6\m@th$}%
	\sbox4{$#7\dabar@\m@th$}%
	\dimen@=\wd0 %
	\ifdim\wd2 >\dimen@
	\dimen@=\wd2 %   
	\fi
	\count@=2 %
	\def\da@bars{\dabar@\dabar@}%
	\@whiledim\count@\wd4<\dimen@\do{%
		\advance\count@\@ne
		\expandafter\def\expandafter\da@bars\expandafter{%
			\da@bars%
			\dabar@%
		}%
	}%  
	\mathrel{#3}%
	\mathrel{%   
		\mathop{\da@bars}\limits
		\ifx\\#1\\%
		\else%
		_{\copy0}%
		\fi
		\ifx\\#2\\%
		\else%
%		^{\copy2}%
		  \ifx#7\scriptstyle%
		    ^{\smash{\substack{\\[0pt]\copy2}}} %adjust vertical positioning of superscript
		  \else%		
		     ^{\smash{\substack{\\[2pt]\copy2}}} %adjust vertical positioning of superscript
		  \fi
		\fi
	}%   
	\mathrel{#4}%
\ifx#7\scriptstyle{}\else{\!\!}\fi% remove some space after the arrow
}
\makeatother

% \cupdot and bigcupdot. See http://tex.stackexchange.com/questions/110976/is-a-cupdot-symbol-available-in-amsmath
\makeatletter
\providecommand*{\cupdot}{%
  \mathbin{%
    \mathpalette\@cupdot{}%
  }%
}
\newcommand*{\@cupdot}[2]{%
  \ooalign{%
    $\m@th#1\cup$\cr
    \sbox0{$#1\cup$}%
    \dimen@=\ht0 %
    \sbox0{$\m@th#1\cdot$}%
    \advance\dimen@ by -\ht0 %
    \dimen@=.5\dimen@
    \hidewidth\raise\dimen@\box0\hidewidth
  }%
}
\providecommand*{\bigcupdot}{%
  \mathop{%
    \vphantom{\bigcup}%
    \mathpalette\@bigcupdot{}%
  }%
}
\newcommand*{\@bigcupdot}[2]{%
  \ooalign{%
    $\m@th#1\bigcup$\cr
    \sbox0{$#1\bigcup$}%
    \dimen@=\ht0 %
    \advance\dimen@ by -\dp0 %
    \sbox0{\scalebox{2}{$\m@th#1\cdot$}}%
    \advance\dimen@ by -\ht0 %
    \dimen@=.5\dimen@
    \hidewidth\raise\dimen@\box0\hidewidth
  }%
}
\makeatother

\newcommand{\RR}{\ensuremath{\mathbb{R}}}
\newcommand{\eRR}{\ensuremath{\overline{\mathbb{R}}}} % extended real numbers
\newcommand{\RRnn}{\ensuremath{\mathbb{R}_{\ge 0}}}
\newcommand{\RRgz}{\ensuremath{\mathbb{R}_{> 0}}}
\newcommand{\RRnp}{\ensuremath{\mathbb{R}_{\le 0}}}
\newcommand{\RRlz}{\ensuremath{\mathbb{R}_{< 0}}}
\newcommand{\RRinfty}{\ensuremath{\mathbb{R}^\infty}}
\newcommand{\NN}{\ensuremath{\mathbb{N}}}
\newcommand{\NNgz}{\ensuremath{\mathbb{N}_{> 0}}}

\newcommand{\powerset}[1]{\ensuremath{2^{#1}}}

\newcommand{\diff}{\ensuremath{\,\mathrm{d}}}

\newcommand{\equalityMarker}{\ensuremath{\ast}}
\newcommand{\eqInductionHypothesis}{\ensuremath{\overset{\mathclap{\mathit{IH}}}{=}}}
\newcommand{\geInductionHypothesis}{\ensuremath{\overset{\mathclap{\mathit{IH}}}{\ge}}}
\newcommand{\leInductionHypothesis}{\ensuremath{\overset{\mathclap{\mathit{IH}}}{\le}}}

\newcommand{\true}{\ensuremath{\mathit{true}}}
\newcommand{\false}{\ensuremath{\mathit{false}}}

%%%%%%%%%%%%%% Probability Theory

\newcommand{\dist}{\ensuremath{\mu}}
\newcommand{\Dist}[1]{\ensuremath{\mathit{Dist}(#1)}}

\newcommand{\SampleSpace}{\ensuremath{\Omega}}
\newcommand{\Events}{\ensuremath{\mathcal{E}}}
\newcommand{\event}{\ensuremath{{E}}}

\newcommand{\rv}{\ensuremath{X}}

%%%%%%%%%%%% experiments
\newcommand{\tool}[1]{\texttt{#1}\xspace}
\newcommand{\prism}{\tool{PRISM}}
\newcommand{\storm}{\tool{Storm}}
\newcommand{\mrmc}{\tool{MRMC}}

\newcommand{\TO}{{\centering TO}}

% You can define your own commands. This is useful for shorthands, for example
\newcommand{\Exp}{\mathbb{E}}
\newcommand{\ExpUpper}{\mathbb{\overline{E}}}
\newcommand{\ExpLower}{\mathbb{\underline{E}}}
\newcommand{\CC}{\mathbb{C}}
\newcommand{\MC}{\mathcal{D}} %Markov Chain
\newcommand{\IMC}{\mathcal{I}} %Interval MC
\newcommand{\MDP}{\mathcal{M}} %MDP
\newcommand{\AIMC}{\mathcal{A}} %AIMC
\newcommand{\tT}{\mathcal{T}} %Transition Matrix
\newcommand{\eE}{\mathcal{E}} %Vector of EVTs
\newcommand{\gG}{\mathcal{G}} %graph
\newcommand{\gS}{\mathcal{S}} %sinks
\newcommand{\nN}{\mathcal{N}}
\newcommand{\PF}{\mathbf{P}} % Probability function
\newcommand{\PFupper}{\mathbf{\overline{P}}}
\newcommand{\PFlower}{\mathbf{\underline{P}}}
\newcommand{\Prupper}{\overline{Pr}}
\newcommand{\Prlower}{\underline{Pr}}
\newcommand{\Initupper}{\mathbf{\overline{\iota}}_{init}}
\newcommand{\Initlower}{\mathbf{\underline{\iota}}_{init}}
\newcommand{\Thetaupper}{\overline{\theta}}
\newcommand{\Thetalower}{\underline{\theta}}
\newcommand{\cCircle}{\bigcirc} % Probability function
\newcommand{\Path}{Paths^{\MC}_{inf}} % Probability function
\newcommand{\lnorm}[1]{\ell_{#1}\text{-norm}}

\newcommand{\T}{\mathsf{T}}   
\renewcommand{\vec}[1]{\bm{#1}}        % redefine \vec to produce boldface instead of adding an arrow
\newcommand{\mat}[1]{\bm{#1}} 
\newcommand{\eye}{\bm{I}}
\newcommand{\Strin}{S_{tr\cap in}}
\newcommand{\until}{\mathbin{\mathsf{U}}}